% Cote 119: offprints, reports, courses, CV and letters (typed and printed part): archive cote 119, site page 142 (PDF page 143).
% typescript. Draft compared: none.
% AI-transcribed by Claude Opus 5.5 (claude-opus-5-5, adaptive thinking, effort high), Message Batches API, batch r119_full (msgbatch_01DwZNm5Y7YsKBWyAN4sFs8w), 2026-11-13; not yet checked by eye.
\scanpage{119/142}
En dehors de cet enseignement,  j'ai assuré et assure des directions de
recherche au niveau de thèses de 3$^{\text{e}}$ cycle ou de doctorat d'état,  sur des
sujets assez divers, dont le fil conducteur commun (à une exception près)
est la topologie des surfaces\xout{,} et des courbes tracées sur les surfaces,
% hand: ``au premier plan'' circled, with a line to a caret after ``avec''
avec\inserted{au premier plan} l'aspect combinatoire des situations envisagées au premier plan.

Ma réflexion personnelle,  en dehors de celle suscitée par mon rôle de
``direction'' de recherche, a porté principalement sur les sujets suivants:

a)  \unsure{Etude|Étude} des actions de groupes de Galois profinis (et notamment de
celui de $\overline{Q}$ sur Q) sur des groupes fondamentaux profinis non abéliens,  tels
les groupes fondamentaux géométriques de courbes algébriques,  et développe-%
ment d'un vaste programme de ``géométrie algébrique anabélienne'', dont un des
éléments cruciaux consiste en des analogues non abéliens des conjectures
de Tate en cohomologie \hand{$\ell$}-adique. A l'horizon se trouve,  entre autres, une
description purement algébrique des schémas de type fini sur la base absolue
Spec($\underline{Z}$) (plus précisément,  de la catégorie des fractions déduite de la caté
gorie formée de ces schémas,  en rendant inversibles les ``homéomorphismes
universels'' i.e. les morphismes qui induisent une équivalence entre les \xout{\unsure{sites}}
topos étales associés) en termes de groupes profinis,  par une description
qui en principe serait purement algébrique (sans référence explicite à des
systèmes d'équations algébriques sur $\underline{Z}$ ...),  et impliquerait une descripti-%
on purement algébrique du groupe Gal$_{\overline{Q}/Q}$ ,  qui dans cette théorie joue le
rôle de la ``base absolue''.  La théorie (en large partie heuristique \unsure{encore})
amène à étudier les classiques groupes de Teichmüller $T_{g,n}$  tant discrets
que profinis,  via les compactifications de Mumford-Deligne des topos modu-%
laires correspondants $M_{g,n}$ ,  et les stratifications naturelles de ceux-ci.
Celles-ci donnent lieu à des dévissages canoniques des groupes, ou mieux,
des \uline{groupoides} de Teichmüller,  qui d'une part permettent d'en donner une
écriture par ``générateurs et relations'' adaptée à la structure stratifiée
(``tour de Teichmüller''),  et qui d'autre part``ont une signification arithmé\xout{-}%
tique'' i.e. se reflètent au niveau des groupes profinis correspondants,  et
sont préservées par l'action naturelle de \hand{$\Gamma$} = Gal$_{\overline{Q}/Q}$ .  Un point important
acquis dès à présent est que ce dernier groupe se réalise fidèlement comme
un sous-groupe du groupe des au\xout{x}tomorphismes extérieurs du groupe fondamen-%
tal profini de la droite projective sur $\overline{Q}$ privée de trois points, et que ce
groupe d'au\xout{t}tomorphismes, \xout{peut se décrire} ou plutôt le sous-groupe qui préserv\unsure{e}
la ``structure à lacets'' naturelle de ce groupe profini,  peut se décrire en
termes de ``paramètres\xout{''} profinis''naturels.  Un des problèmes essentiels est
de parvenir à une caractérisation purement algébrique de ce sous-groupe
Gal$_{\overline{Q}/Q}$ ,  et une motivation principale pour l'étude de la ``tour de Teichmül-%
ler'' est de parvenir à une caractérisation heuristique plausible.

b)  En marge de l'étude de la ``tour de Teichmüller'', \deleted{établissement d'une}\inserted{réflexions sur}
théorie générale de ``dévissage'' des structures stratifiées, dans deux con-%
textes qui se complètent mutuellement:  celui des topos,  et celui des ``es-%
paces topologiques M-modérés'',  pour une notion de ``modération'' M convenable
En même temps, réflexions sur les fondements d'une ``topologie modérée'', plus
appropriée à l'étude des espaces d'homéomorphismes ou de plongements,  par
exemple,  et de la relation d'isotopie,  que la notion d'espace topologique,
inadéquate également pour une théorie des structures stratifiées. L'impor-%
tance de ces dernières peut se mesurer à cette observation\hand{,} que la plupart
des ``espaces'' rencontrés,  et plus particulièrement les espaces ou topos
``modulaires'' de tout poil,  sont munis de stratifications naturelles, faisant
partie intégrante de leur structure tant soit peu fine.

c)   Réflexions sur les fondements de l'algèbre homotopique,  dans l'es-%
prit de``l'algèbre cohomologique non commutative de dimension quelconque'',
avec la place centrale à la notion de \hand{$\infty$}-Champ et \hand{$\infty$}-Gr-champ (au sens abso-%
lu,  ou relativisé au dessus d'un topos),  jouant le rôle à la fois de ``mol-%
les'' particulièrement souples pour exprimer des types d'homotopie,  et des
``coéfficients'' les plus généraux pour une ``cohomologie non commutative''.

\mspage{142}
